\section{Simulator validation}
\label{app:simvalidation}

This appendix records the simulator-validation axes and their current
status from \texttt{outputs/tables/bench\_simulator\_validation.csv}.
Four of five axes are \emph{verified} on the chromatin-aware simulator
(\texttt{simulate\_dataset\_chromatin\_aware}); the remaining axis is
\emph{preliminary} pending length-mixture re-fitting against real
cfDNA data (cf.\ Sec.~\ref{sec:discussion}).

\begin{center}
\small
\begin{tabular}{lll}
\toprule
Axis & Measured value & Reference target / status \\
\midrule
Fragment-length primary mode & $167.5$~bp & $\sim 167$~bp \citep{snyder2016}; verified \\
Fragment-length $320$--$350$~bp peak height & $0.0029$ per bp & $\sim 0.005$ target; preliminary (re-fit) \\
Helical-pitch periodicity peak (lag 8--13 bp) & $0.86$ autocorr & non-trivially positive; verified \\
Top-4 4-mer fraction at $5'$ cuts & $0.156$ & $\sim 0.20$ \citep{zhou2022}; verified \\
Methylation-conditioned GC effect size & $0.090$ & non-zero; verified \\
\bottomrule
\end{tabular}
\end{center}

\subsection*{Fragment-length distribution}
The chromatin-aware simulator reproduces a dominant mononucleosomal mode
at $167.5$~bp matching the canonical $167$~bp from \citet{snyder2016}.
The secondary $320$--$350$~bp peak height measures $0.0029$ per bp,
below the published target $\sim 0.005$ because the released 3-mode
mixture parameters are spec-derived (mixture weights $\sim 0.2$ on
$\mathcal{N}(332, 30^2)$ gives $\sim 0.003$ by construction). Closing
this gap requires re-fitting the length-mixture parameters on real
Liu 2024 fragments; this is tracked as an open follow-up.

\subsection*{End-motif frequencies}
The chromatin-aware simulator's top-4 4-mer fraction at $5'$ cut sites
reaches $0.156$, close to the $\sim 0.20$ canonical fraction from
\citet{zhou2022}. The boosted-logit baseline parameter controls this
fraction directly; full motif-context tuning is available.

\subsection*{Periodicity spectrum}
Cut-site autocorrelation in lag $8$--$13$~bp exhibits a peak of $0.86$
under the chromatin-aware path (versus $0.84$ under the legacy harness,
both verified). The peak comes from \texttt{cut\_site\_density\_linker}
which implements the App.~E corrected formula $p_{\mathrm{cut}}(x)
\propto \exp(-d_{\mathrm{linker}}/\sigma_{\mathrm{link}})\cdot
[1+a\cos(2\pi d_{\mathrm{center}}/10.4)]_+$ -- peaks fall at midpoints
between consecutive nucleosomes, not at their centers.

\subsection*{Methylation-conditioned cut bias}
A methylation-conditioned GC effect size of $0.090$ is measured at
the canonical seed. This is the signal exploited by FinaleMe-style and
HAVI-Methyl predictors and is the primary reason the model can recover
$\beta$ from fragmentomic features alone in the deep-coverage regime.
